\chapter{नाभिकरागी प्रतिस्थापन}\label{ch:b1:nucleophilic-substitution}

\cip{S}-2-ब्रोमोब्यूटेन की एक बोतल लीजिए, एक अकेला \omterm{def:b1:stereochemistry-in-depth:enantiomers}{प्रतिबिंबरूपी} जो
ध्रुवित प्रकाश को घुमाता है। ध्रुवीय \omterm{def:b1:intermolecular-forces-solvents:solvent-classes}{अप्रोटिक विलायक} में सोडियम हाइड्रॉक्साइड के साथ गरम करने पर
यह ब्यूटेन-2-ऑल देता है जो भी एक अकेला \omterm{def:b1:stereochemistry-in-depth:enantiomers}{प्रतिबिंबरूपी} है --- पर दूसरा वाला,
\cip{R}: अभिक्रिया ने अणु को आँधी में उलटे छाते की तरह भीतर से बाहर
पलट दिया है। इसके बजाय गुनगुने पानी में छोड़ देने पर, वही ब्रोमाइड ऐसा
ब्यूटेन-2-ऑल देता है जो आंशिक रूप से रेसेमिक है। एक ही समग्र परिवर्तन, Br परमाणु का
स्थान लेता एक OH समूह, दो भिन्न मार्गों से होता है। यह अध्याय
दोनों क्रियाविधियों को उनके \omterm{def:b1:rate-laws:order}{दर नियमों} से स्थापित करता है, उनके
त्रिविम रसायन की व्याख्या करता है, और वे नियम देता है जो उनके बीच निर्णय करते हैं --- यह
गतिकी और इलेक्ट्रॉनिक प्रभावों वाले अध्यायों के उपकरणों का
पहला पूर्ण उपयोग है।

\begin{recall}
\omterm{def:b1:rate-laws:order}{दर नियम}, \omterm{def:b1:elementary-steps:elementary-step}{मूल चरण}, \omterm{def:b1:elementary-steps:rds}{दर-निर्धारक चरण} और
\omterm{def:b1:elementary-steps:qssa}{स्थायी-अवस्था सन्निकटन} \cref{ch:b1:rate-laws,ch:b1:elementary-steps} में हैं;
\omterm{def:b1:stereochemistry-in-depth:isomers}{विन्यास}, क्रैम चित्र और \omterm{def:b1:stereochemistry-in-depth:optical-activity}{प्रकाशिक सक्रियता}
\cref{ch:b1:stereochemistry-in-depth} में; \omterm{def:b1:electronic-effects:nucleophile}{नाभिकरागी}, \omterm{def:b1:electronic-effects:nucleophile}{इलेक्ट्रॉनरागी},
\omterm{def:b1:electronic-effects:nucleophile}{निर्गामी समूह} और \omterm{def:b1:electronic-effects:intermediates}{कार्बधनायन} \cref{ch:b1:electronic-effects} में;
प्रोटिक और \omterm{def:b1:intermolecular-forces-solvents:solvent-classes}{अप्रोटिक विलायक} \cref{ch:b1:intermolecular-forces-solvents} में।
\end{recall}

\section{प्रतिस्थापन और उसके भागीदार}

\begin{definition}[नाभिकरागी प्रतिस्थापन, क्रियाधार]\label{def:b1:nucleophilic-substitution:substitution}
\emph{नाभिकरागी प्रतिस्थापन}\index{नाभिकरागी प्रतिस्थापन}
किसी संतृप्त कार्बन पर \omterm{def:b1:electronic-effects:nucleophile}{निर्गामी समूह} Y का \omterm{def:b1:electronic-effects:nucleophile}{नाभिकरागी} Nu द्वारा
प्रतिस्थापन है: $\mathrm{Nu^-} + \mathrm{R{-}Y} \to \mathrm{R{-}Nu} +
\mathrm{Y^-}$। यौगिक R--Y \emph{क्रियाधार}\index{क्रियाधार} है;
Y वहन करने वाला कार्बन \omterm{def:b1:electronic-effects:nucleophile}{इलेक्ट्रॉनरागी} होता है क्योंकि C--Y आबंध
Y की ओर ध्रुवित है।
\end{definition}

हैलोऐल्केन विशिष्ट \omterm{def:b1:nucleophilic-substitution:substitution}{क्रियाधार} हैं, और उनमें \ce{Cl-}, \ce{Br-} या
\ce{I-} \omterm{def:b1:electronic-effects:nucleophile}{निर्गामी समूह} की भूमिका निभाते हैं। हाइड्रॉक्साइड ऐल्कोहॉल देता है; ऐल्कॉक्साइड
ईथर देते हैं; सायनाइड नाइट्राइल देता है (एक अधिक कार्बन); अमोनिया और ऐमीन
ऐमीन देते हैं; आयोडाइड अन्य हैलाइडों से विनिमय करता है। विलायक के रूप में प्रयुक्त
जल और ऐल्कोहॉल भी \omterm{def:b1:electronic-effects:nucleophile}{नाभिकरागी} की भूमिका निभा सकते हैं: तब प्रतिस्थापन
एक \omterm{def:b1:nucleophilic-substitution:sn1}{विलायक-अपघटन} होता है।

\section{SN2 क्रियाविधि}

जब ब्रोमोमेथेन जल में हाइड्रॉक्साइड से अभिक्रिया करता है, तो किसी भी
सांद्रता को दुगुना करने से दर दुगुनी हो जाती है: $v = k[\ce{CH3Br}][\ce{OH-}]$, एक ऐसी अभिक्रिया
जिसकी हर अभिकारक में कोटि एक है।

\begin{definition}[SN2 क्रियाविधि, वाल्डेन प्रतीपन]\label{def:b1:nucleophilic-substitution:sn2}
\emph{SN2 क्रियाविधि}\index{SN2 क्रियाविधि} (प्रतिस्थापन,
\omterm{def:b1:electronic-effects:nucleophile}{नाभिकरागी}, द्विअणुक) एक अकेला \omterm{def:b1:elementary-steps:elementary-step}{मूल चरण} है जिसमें
\omterm{def:b1:electronic-effects:nucleophile}{नाभिकरागी} \omterm{def:b1:electronic-effects:nucleophile}{निर्गामी समूह} के विपरीत ओर से कार्बन पर आक्रमण करता है,
जबकि \omterm{def:b1:electronic-effects:nucleophile}{निर्गामी समूह} निकल जाता है। कार्बन के तीन अन्य प्रतिस्थापी
उलट जाते हैं, जैसे भीतर से बाहर पलटा छाता: \emph{वाल्डेन
प्रतीपन}\index{वाल्डेन प्रतीपन}।
\end{definition}

\begin{omfigure}
\setchemfig{atom sep=2.1em}
\schemestart
\chemfig{H@{nu}O^{-}}
\arrow{0}[,0.35]
\chemfig{@{c}C(-[:150]H_3C)(<[:210]H)(<:[:250]C_2H_5)-[@{cb}]@{br}Br}
\arrow{->}
\chemfig{[:0]HO-C(-[:30]CH_3)(<[:-30]H)(<:[:-70]C_2H_5)}
\+
\chemfig{Br^{-}}
\schemestop
\chemmove{\draw[->, thick, shorten <=2pt, shorten >=3pt](nu) .. controls +(15:5mm) and +(175:5mm) .. (c);
          \draw[->, thick, shorten <=2pt](cb) .. controls +(90:6mm) and +(90:6mm) .. (br);}

\medskip
\begin{tikzpicture}[font=\small]
  \node[draw, dashed, rounded corners, inner sep=6pt] at (0,0)
    {$\left[\vcenter{\hbox{\ce{HO}$\cdots$\chemfig{C(-[2]CH_3)(-[:-30]H)(-[:210]C_2H_5)}$\cdots$\ce{Br}}}\right]^{-}$};
  \node[below] at (0,-1.1) {संक्रमण अवस्था: C के चारों ओर पाँच समूह, तीन एक तल में};
\end{tikzpicture}

{\small \cip{S}-2-ब्रोमोब्यूटेन के साथ हाइड्रॉक्साइड की SN2 अभिक्रिया।
\omterm{def:b1:electronic-effects:nucleophile}{नाभिकरागी} ब्रोमीन के विपरीत फलक पर आक्रमण करता है (\omterm{def:b1:electronic-effects:curly-arrow}{वक्र तीर}); \omterm{def:b1:elementary-steps:energy-profile}{संक्रमण
अवस्था} में O--C आबंध बन रहा होता है जबकि C--Br आबंध टूटता है;
उत्पाद \cip{R}-ब्यूटेन-2-ऑल है: तीनों अन्य समूह पलटकर
दूसरी ओर आ गए हैं।}
\end{omfigure}

\begin{proposition}[SN2 का दर नियम]\label{prop:b1:nucleophilic-substitution:sn2-law}
SN2 अभिक्रिया के लिए $v = k[\mathrm{RY}][\mathrm{Nu}]$।
\end{proposition}

\begin{proof}
क्रियाविधि एक अकेला द्विअणुक \omterm{def:b1:elementary-steps:elementary-step}{मूल चरण} है; \omterm{def:b1:elementary-steps:elementary-step}{मूल चरणों} के लिए
द्रव्य-अनुपाती क्रिया के नियम (\cref{ch:b1:elementary-steps}) से उसकी दर
दोनों भागीदारों की सांद्रताओं के गुणनफल के समानुपाती है।
\end{proof}

\begin{proposition}[SN2 का त्रिविम रसायन]\label{prop:b1:nucleophilic-substitution:sn2-inversion}
त्रिविमजनक कार्बन पर SN2 अभिक्रिया एक अकेला \omterm{def:b1:stereochemistry-in-depth:isomers}{त्रिविम समावयवी} देती है, जिसमें
तीन अपरिवर्तित समूहों की व्यवस्था प्रतीपित होती है: इस प्रकार की
अभिक्रिया त्रिविमविशिष्ट है।
\end{proposition}

\begin{proof}
\omterm{def:b1:electronic-effects:nucleophile}{नाभिकरागी} केवल कार्बन की रिक्त ओर तक पहुँच सकता है, जो
\omterm{def:b1:electronic-effects:nucleophile}{निर्गामी समूह} के विपरीत है; तीनों अन्य समूह \omterm{def:b1:elementary-steps:energy-profile}{संक्रमण अवस्था} में एक तल से
गुज़रते हैं और दूसरी ओर पहुँच जाते हैं। नया आबंध वहाँ संकेत करता है जहाँ
पुराना नहीं करता था: \omterm{def:b1:nucleophilic-substitution:substitution}{क्रियाधार} का हर अणु अपने समूहों की दर्पण-प्रतिबिंब
व्यवस्था देता है। वर्णक प्रायः \cip{S} से \cip{R} में या उलटा बदल जाता है,
पर हमेशा नहीं, क्योंकि वह \omterm{def:b1:electronic-effects:nucleophile}{नाभिकरागी} और \omterm{def:b1:electronic-effects:nucleophile}{निर्गामी समूह} के क्रम पर
निर्भर करता है।
\end{proof}

\begin{definition}[त्रिविमविशिष्ट अभिक्रिया]\label{def:b1:nucleophilic-substitution:stereospecific}
कोई अभिक्रिया \emph{त्रिविमविशिष्ट अभिक्रिया}\index{त्रिविमविशिष्ट अभिक्रिया}
है जब त्रिविमसमावयवी \omterm{def:b1:nucleophilic-substitution:substitution}{क्रियाधार} त्रिविमसमावयवी रूप से
भिन्न उत्पाद दें, हर \omterm{def:b1:nucleophilic-substitution:substitution}{क्रियाधार} अपना अलग उत्पाद दे।
\end{definition}

\section{SN1 क्रियाविधि}

2-ब्रोमो-2-मेथिलप्रोपेन जल से ऐसी दर पर अभिक्रिया करता है जो
\omterm{def:b1:electronic-effects:nucleophile}{नाभिकरागी} की सांद्रता पर निर्भर नहीं करती: $v = k[(\ce{CH3})_3\ce{CBr}]$।
दर का निर्णय किसी ऐसे चरण को करना चाहिए जिसमें \omterm{def:b1:electronic-effects:nucleophile}{नाभिकरागी} भाग नहीं लेता।

\begin{definition}[SN1 क्रियाविधि, रेसेमीकरण, विलायक-अपघटन]\label{def:b1:nucleophilic-substitution:sn1}
\emph{SN1 क्रियाविधि}\index{SN1 क्रियाविधि} (प्रतिस्थापन,
\omterm{def:b1:electronic-effects:nucleophile}{नाभिकरागी}, एकाणुक) के दो चरण होते हैं: C--Y आबंध का मंद \omterm{def:b1:electronic-effects:cleavage}{विषमांश विदलन},
जो एक \omterm{def:b1:electronic-effects:intermediates}{कार्बधनायन} देता है, फिर \omterm{def:b1:electronic-effects:intermediates}{कार्बधनायन} पर \omterm{def:b1:electronic-effects:nucleophile}{नाभिकरागी} का
तीव्र योग। जो अभिक्रिया एक अकेले
\omterm{def:b1:stereochemistry-in-depth:enantiomers}{प्रतिबिंबरूपी} को दोनों के मिश्रण में बदल दे, वह \emph{रेसेमीकरण}\index{रेसेमीकरण} है;
जिस प्रतिस्थापन में विलायक ही \omterm{def:b1:electronic-effects:nucleophile}{नाभिकरागी} हो, वह
\emph{विलायक-अपघटन}\index{विलायक-अपघटन} है।
\end{definition}

\begin{omfigure}
\setchemfig{atom sep=2.1em}
\schemestart
\chemfig{C(-[2]CH_3)(-[4]H_3C)(-[6]CH_3)-[@{cb2}]@{br2}Br}
\arrow{->[मंद]}
\chemfig{C^{+}(-[2]CH_3)(-[:210]H_3C)(-[:-30]CH_3)}
\+
\chemfig{Br^{-}}
\schemestop

\smallskip
\schemestart
\chemfig{C^{+}(-[2]CH_3)(-[:210]H_3C)(-[:-30]CH_3)}
\arrow{->[तीव्र][\ce{H2O}]}
\chemfig{C(-[2]CH_3)(-[4]H_3C)(-[6]CH_3)-OH_2^{+}}
\arrow{->[\ce{-H+}]}
\chemfig{C(-[2]CH_3)(-[4]H_3C)(-[6]CH_3)-OH}
\schemestop
\chemmove{\draw[->, thick, shorten <=2pt](cb2) .. controls +(90:6mm) and +(90:6mm) .. (br2);}

\medskip
\begin{tikzpicture}[font=\small]
  \fill[omProp!15, draw=omProp] (0,0.2) ellipse (0.2 and 0.85);
  \fill[omProp!15, draw=omProp] (0,-0.2) ellipse (0.2 and 0.85);
  \draw[thick] (0,0) -- (-1.1,-0.2) node[left] {\ce{CH3}};
  \draw[thick] (0,0) -- (0.95,-0.45) node[right] {\ce{C2H5}};
  \draw[thick] (0,0) -- (0.4,0.5) node[above right] {H};
  \node[fill=white, inner sep=1pt] at (0,0) {\ce{C+}};
  \draw[->, thick, omDef] (-0.9,1.6) node[left] {\ce{H2O}} -- (-0.15,0.95);
  \draw[->, thick, omDef] (0.9,-1.6) node[right] {\ce{H2O}} -- (0.15,-0.95);
  \node[align=center] at (4.6,0) {किसी भी फलक पर आक्रमण:\\ दोनों विन्यास,\\ रेसेमिक, यदि फलकों के बीच\\ और कोई भेद न हो};
\end{tikzpicture}

{\small जल में 2-ब्रोमो-2-मेथिलप्रोपेन का SN1 \omterm{def:b1:nucleophilic-substitution:sn1}{विलायक-अपघटन}: तृतीयक \omterm{def:b1:electronic-effects:intermediates}{कार्बधनायन}
तक मंद आयनन, फिर जल द्वारा ग्रहण और एक प्रोटॉन की हानि।
नीचे: एक \omterm{def:b1:stereochemistry-in-depth:stereocentre}{किरैल} द्वितीयक \omterm{def:b1:electronic-effects:intermediates}{कार्बधनायन} (2-ब्रोमोब्यूटेन से) समतलीय है,
और जल किसी भी फलक पर समान प्रायिकता से जुड़ता है।}
\end{omfigure}

\begin{proposition}[SN1 का दर नियम]\label{prop:b1:nucleophilic-substitution:sn1-law}
\omterm{def:b1:nucleophilic-substitution:sn1}{SN1 क्रियाविधि} के लिए, यदि \omterm{def:b1:electronic-effects:intermediates}{कार्बधनायन} \omterm{def:b1:nucleophilic-substitution:substitution}{क्रियाधार} में लौटने की तुलना में
कहीं अधिक तेज़ी से ग्रहण कर लिया जाता है, तो $v = k_1[\mathrm{RY}]$, जो
\omterm{def:b1:electronic-effects:nucleophile}{नाभिकरागी} से स्वतंत्र है।
\end{proposition}

\begin{proof}
चरण: $\mathrm{RY} \rightleftharpoons \mathrm{R^+} + \mathrm{Y^-}$ ($k_1$,
$k_{-1}$), $\mathrm{R^+} + \mathrm{Nu} \to \mathrm{RNu}$ ($k_2$)।
\omterm{def:b1:electronic-effects:intermediates}{कार्बधनायन} एक \omterm{def:b1:electronic-effects:intermediates}{क्रियाशील मध्यवर्ती} है: \omterm{def:b1:elementary-steps:qssa}{स्थायी-अवस्था
सन्निकटन} (\cref{ch:b1:elementary-steps}) से $k_1[\mathrm{RY}] =
k_{-1}[\mathrm{R^+}][\mathrm{Y^-}] + k_2[\mathrm{R^+}][\mathrm{Nu}]$, इसलिए
$v = k_2[\mathrm{R^+}][\mathrm{Nu}] = \dfrac{k_1k_2[\mathrm{RY}][\mathrm{Nu}]}
{k_{-1}[\mathrm{Y^-}] + k_2[\mathrm{Nu}]}$। जब $k_2[\mathrm{Nu}] \gg
k_{-1}[\mathrm{Y^-}]$ हो, तब यह $k_1[\mathrm{RY}]$ रह जाता है: पहला चरण
दर-निर्धारक है।
\end{proof}

\begin{proposition}[SN1 का त्रिविम रसायन]\label{prop:b1:nucleophilic-substitution:sn1-stereo}
त्रिविमजनक कार्बन पर SN1 अभिक्रिया दोनों \omterm{def:b1:stereochemistry-in-depth:isomers}{विन्यास} देती है:
अभिक्रिया त्रिविमविशिष्ट नहीं है, और जब \omterm{def:b1:electronic-effects:intermediates}{कार्बधनायन} के दोनों फलकों में कोई भेद
न हो, तो \omterm{def:b1:nucleophilic-substitution:sn1}{रेसेमीकरण} की ओर ले जाती है।
\end{proposition}

\begin{proof}
\omterm{def:b1:electronic-effects:intermediates}{कार्बधनायन} समतलीय है (\cref{ch:b1:electronic-effects}): उसके दोनों
फलक दर्पण-प्रतिबिंब हैं। यदि ग्रहण से पहले \omterm{def:b1:electronic-effects:nucleophile}{निर्गामी समूह} दूर जा चुका हो,
तो \omterm{def:b1:electronic-effects:nucleophile}{नाभिकरागी} किसी भी फलक पर समान दर से जुड़ता है, जिससे
दोनों \omterm{def:b1:stereochemistry-in-depth:enantiomers}{प्रतिबिंबरूपी} समान मात्रा में मिलते हैं। व्यवहार में निर्गामी आयन कुछ समय तक
एक फलक को ढके रहता है, और प्रायः प्रतीपन का थोड़ा आधिक्य
प्रेक्षित होता है।
\end{proof}

\begin{omfigure}
\begin{tikzpicture}
\begin{axis}[
  width=0.8\linewidth, height=5.0cm, xmin=0, xmax=10, ymin=-30, ymax=110,
  axis x line=bottom, axis y line=left, xlabel={अभिक्रिया निर्देशांक},
  ylabel={ऊर्जा}, xtick=\empty, ytick=\empty,
  label style={font=\small}, legend style={font=\scriptsize, draw=none, at={(0.98,0.98)}, anchor=north east}]
\addplot[omDef, very thick] table[x=x, y=E] {figdata/out/bachelor-1-energy-profiles-sn2.dat};
\addplot[omProp, very thick, dashed] table[x=x, y=E] {figdata/out/bachelor-1-energy-profiles-sn1.dat};
\legend{SN2: एक चरण, SN1: दो चरण}
\node[font=\scriptsize, anchor=north] at (axis cs:5,68) {कार्बधनायन};
\node[font=\scriptsize, anchor=north] at (axis cs:0.6,-3) {RY};
\node[font=\scriptsize, anchor=south] at (axis cs:9.6,-18) {RNu};
\end{axis}
\end{tikzpicture}

{\small योजनात्मक \omterm{def:b1:elementary-steps:energy-profile}{ऊर्जा प्रोफ़ाइल}। SN2: एक अकेली \omterm{def:b1:elementary-steps:energy-profile}{संक्रमण अवस्था}, कार्बन के चारों ओर पाँच
समूह। SN1: पहला, ऊँचा अवरोध (आयनन, जो
\omterm{def:b1:elementary-steps:rds}{दर-निर्धारक चरण} है), एक उथले गर्त में \omterm{def:b1:electronic-effects:intermediates}{कार्बधनायन}, फिर उसके ग्रहण तक एक छोटा
अवरोध।}
\end{omfigure}

\section{क्रियाविधि का निर्णय किससे होता है}

\begin{proposition}[क्रियाधार]\label{prop:b1:nucleophilic-substitution:substrate}
मेथिल और प्राथमिक \omterm{def:b1:nucleophilic-substitution:substitution}{क्रियाधार} SN2 से अभिक्रिया करते हैं; तृतीयक \omterm{def:b1:nucleophilic-substitution:substitution}{क्रियाधार} SN1 से;
द्वितीयक \omterm{def:b1:nucleophilic-substitution:substitution}{क्रियाधार} किसी से भी, अन्य कारकों के अनुसार। ऐलिलिक
और बेंज़िलिक \omterm{def:b1:nucleophilic-substitution:substitution}{क्रियाधार} दोनों से तेज़ी से अभिक्रिया करते हैं।
\end{proposition}

\begin{proof}
SN2 को कार्बन के पीछे स्थान चाहिए: हाइड्रोजन के स्थान पर हर ऐल्किल समूह
पाँच समूहों वाली \omterm{def:b1:elementary-steps:energy-profile}{संक्रमण अवस्था} में भीड़ बढ़ाता है और
अभिक्रिया को धीमा करता है; तीन समूह उसे रोक देते हैं। SN1 को स्थायी \omterm{def:b1:electronic-effects:intermediates}{कार्बधनायन} चाहिए:
स्थायित्व का क्रम (\cref{prop:b1:electronic-effects:carbocation-order})
आयनन को तृतीयक के लिए तीव्र, द्वितीयक के लिए संभव, और प्राथमिक तथा मेथिल
\omterm{def:b1:electronic-effects:intermediates}{कार्बधनायनों} के लिए बहुत धीमा बनाता है। ऐलिलिक और बेंज़िलिक \omterm{def:b1:electronic-effects:intermediates}{कार्बधनायन}
अनुनाद द्वारा स्थायी होते हैं, और वही $\pi$ निकाय SN2
\omterm{def:b1:elementary-steps:energy-profile}{संक्रमण अवस्था} को भी स्थायी बनाता है।
\end{proof}

\begin{proposition}[निर्गामी समूह]\label{prop:b1:nucleophilic-substitution:leaving-group}
\omterm{def:b1:electronic-effects:nucleophile}{निर्गामी समूह} जितना अच्छा, अर्थात् वह जितना \omterm{def:b1:predominant-reaction:strong-acid}{दुर्बल क्षारक} बनता है,
दोनों क्रियाविधियाँ उतनी ही तीव्र होती हैं: $\ce{I-} > \ce{Br-} > \ce{Cl-} \gg \ce{F-}$;
\ce{OH-} और \ce{RO-} नहीं निकलते।
\end{proposition}

\begin{proof}
दोनों में C--Y आबंध \omterm{def:b1:elementary-steps:rds}{दर-निर्धारक चरण} में या उससे पहले टूटता है,
और \omterm{def:b1:elementary-steps:energy-profile}{संक्रमण अवस्था} ऋण आवेश का एक भाग Y पर वहन करती है। जो
समूह ऋण आवेश को अच्छी तरह सँभालता है --- किसी
\omterm{def:b1:predominant-reaction:strong-acid}{प्रबल अम्ल} का संयुग्मी क्षारक --- वह उस \omterm{def:b1:elementary-steps:energy-profile}{संक्रमण अवस्था} को नीचे लाता है। \ce{HI}, \ce{HBr} और
\ce{HCl} \omterm{def:b1:predominant-reaction:strong-acid}{प्रबल अम्ल} हैं, \ce{HF} दुर्बल ($\mathrm{p}K_a$ 3.2), जल और
ऐल्कोहॉल अत्यंत दुर्बल: \ce{OH-} और \ce{RO-} \omterm{def:b1:predominant-reaction:strong-acid}{प्रबल क्षारक} और खराब
\omterm{def:b1:electronic-effects:nucleophile}{निर्गामी समूह} हैं।
\end{proof}

\begin{proposition}[विलायक]\label{prop:b1:nucleophilic-substitution:solvent}
ध्रुवीय \omterm{def:b1:intermolecular-forces-solvents:solvent-classes}{प्रोटिक विलायक} (जल, ऐल्कोहॉल) SN1 के अनुकूल हैं; ध्रुवीय \omterm{def:b1:intermolecular-forces-solvents:solvent-classes}{अप्रोटिक विलायक}
(प्रोपेनोन, डाइमेथिल सल्फ़ॉक्साइड, डाइमेथिलफ़ॉर्मैमाइड) ऋणायनी \omterm{def:b1:electronic-effects:nucleophile}{नाभिकरागियों} के साथ
SN2 के अनुकूल हैं।
\end{proposition}

\begin{proof}
SN1 एक उदासीन \omterm{def:b1:nucleophilic-substitution:substitution}{क्रियाधार} से दो आयन बनाता है: ध्रुवीय \omterm{def:b1:intermolecular-forces-solvents:solvent-classes}{प्रोटिक विलायक}
दोनों को स्थायी बनाता है, धनायन को अपने \omterm{def:b1:lewis-resonance-vsepr:lewis}{एकाकी युग्मों} से और ऋणायन को
\omterm{def:b1:intermolecular-forces-solvents:hbond}{हाइड्रोजन आबंधों} से, और आयनन अवरोध को नीचे लाता है। SN2 में
ऋणायनी \omterm{def:b1:electronic-effects:nucleophile}{नाभिकरागी} का आवेश \omterm{def:b1:elementary-steps:energy-profile}{संक्रमण अवस्था} में फैल जाता है; \omterm{def:b1:intermolecular-forces-solvents:solvent-classes}{प्रोटिक
विलायक}, जो छोटे ऋणायन को \omterm{def:b1:intermolecular-forces-solvents:hbond}{हाइड्रोजन आबंधों} से प्रबलता से बाँधता है, उसे
धीमा करता है, जबकि \omterm{def:b1:intermolecular-forces-solvents:solvent-classes}{अप्रोटिक विलायक} ऋणायन को मुक्त और क्रियाशील छोड़ देता है
(\cref{ch:b1:intermolecular-forces-solvents})।
\end{proof}

\begin{method}[क्रियाविधि चुनना]\label{met:b1:nucleophilic-substitution:choose-mechanism}
\begin{enumerate}
  \item \omterm{def:b1:nucleophilic-substitution:substitution}{क्रियाधार}: मेथिल या प्राथमिक $\to$ SN2; तृतीयक $\to$ SN1;
        द्वितीयक $\to$ आगे बढ़िए।
  \item \omterm{def:b1:electronic-effects:nucleophile}{नाभिकरागी}: प्रबल और ऋणायनी ($\ce{OH-}$, $\ce{RO-}$, $\ce{CN-}$,
        $\ce{I-}$) $\to$ SN2; दुर्बल और उदासीन (जल, ऐल्कोहॉल) $\to$ SN1।
  \item विलायक: अप्रोटिक $\to$ SN2; प्रोटिक $\to$ SN1।
  \item त्रिविम रसायन (प्रतीपन या \omterm{def:b1:nucleophilic-substitution:sn1}{रेसेमीकरण}) का पूर्वानुमान कीजिए और जाँचिए
        कि कोई \omterm{def:b1:predominant-reaction:strong-acid}{प्रबल क्षारक} इसके बजाय विलोपन तो नहीं कराता
        (\cref{ch:b1:elimination})।
\end{enumerate}
\end{method}

\begin{omfigure}
\begin{tabular}{lll}
\hline
 & SN2 & SN1 \\
\hline
चरण & एक & दो, एक \omterm{def:b1:electronic-effects:intermediates}{कार्बधनायन} के माध्यम से \\
\omterm{def:b1:rate-laws:order}{दर नियम} & $k[\mathrm{RY}][\mathrm{Nu}]$ & $k[\mathrm{RY}]$ \\
\omterm{def:b1:nucleophilic-substitution:substitution}{क्रियाधार} & मेथिल $>$ प्राथमिक $>$ द्वितीयक & तृतीयक $>$ द्वितीयक \\
\omterm{def:b1:electronic-effects:nucleophile}{नाभिकरागी} & प्रबल, ऋणायनी & दुर्बल, उदासीन (प्रायः विलायक) \\
विलायक & ध्रुवीय अप्रोटिक & ध्रुवीय प्रोटिक \\
त्रिविम रसायन & प्रतीपन (त्रिविमविशिष्ट) & \omterm{def:b1:nucleophilic-substitution:sn1}{रेसेमीकरण} (अधिकांशतः) \\
\hline
\end{tabular}

\smallskip
{\small दोनों क्रियाविधियाँ साथ-साथ।}
\end{omfigure}

\section{अभ्यास}

\begin{exercise}[$\star$]\label{exo:b1:nucleophilic-substitution:1}
इनके उत्पाद लिखिए: सोडियम हाइड्रॉक्साइड के साथ ब्रोमोएथेन; सोडियम एथॉक्साइड के साथ
आयोडोमेथेन; पोटैशियम सायनाइड के साथ 1-क्लोरोप्रोपेन; अमोनिया के साथ
ब्रोमोमेथेन (पहला चरण)। \omterm{def:b1:nucleophilic-substitution:substitution}{क्रियाधार}, \omterm{def:b1:electronic-effects:nucleophile}{नाभिकरागी} और \omterm{def:b1:electronic-effects:nucleophile}{निर्गामी समूह} पहचानिए।
\end{exercise}

\begin{exercise}[$\star$]\label{exo:b1:nucleophilic-substitution:2}
अभिक्रिया $\mathrm{RY} + \mathrm{Nu}$ के लिए, जिसकी दर $v = k[\mathrm{RY}][\mathrm{Nu}]$ है,
यदि $[\mathrm{Nu}]$ तिगुना कर दिया जाए तो \omterm{def:b1:rate-laws:initial-rate}{प्रारंभिक दर} का क्या होता है? यदि दोनों
सांद्रताएँ आधी कर दी जाएँ? यही प्रश्न, यदि $v = k[\mathrm{RY}]$ हो।
\end{exercise}

\begin{exercise}[$\star$]\label{exo:b1:nucleophilic-substitution:3}
क्रियाविधि (SN1 या SN2) का पूर्वानुमान कीजिए: प्रोपेनोन में सोडियम आयोडाइड के साथ
1-ब्रोमोब्यूटेन; जल में 2-ब्रोमो-2-मेथिलप्रोपेन; हाइड्रॉक्साइड के साथ ब्रोमोमेथेन;
बिना कोई क्षारक मिलाए एथेनॉल में 2-ब्रोमोब्यूटेन।
\end{exercise}

\begin{exercise}[$\star$]\label{exo:b1:nucleophilic-substitution:4}
\cip{R}-2-ब्रोमोब्यूटेन के साथ सायनाइड की SN2 अभिक्रिया का उत्पाद
\omterm{def:b1:stereochemistry-in-depth:representations}{क्रैम निरूपण} में खींचिए और उसका वर्णक बताइए।
\end{exercise}

\begin{exercise}[$\star\star$]\label{exo:b1:nucleophilic-substitution:5}
2-ब्रोमो-2-मेथिलप्रोपेन के जल-अपघटन की पूरी \omterm{def:b1:nucleophilic-substitution:sn1}{SN1 क्रियाविधि} \omterm{def:b1:electronic-effects:curly-arrow}{वक्र तीरों} के साथ लिखिए,
प्रोटॉन की हानि सहित। तीन स्पीशीज़ के भाग लेने के बावजूद \omterm{def:b1:rate-laws:order}{दर नियम}
प्रथम कोटि का क्यों है?
\end{exercise}

\begin{exercise}[$\star\star$]\label{exo:b1:nucleophilic-substitution:6}
समझाइए कि 1-ब्रोमो-2,2-डाइमेथिलप्रोपेन, जो एक प्राथमिक ब्रोमाइड है, SN2 से बहुत
धीरे अभिक्रिया क्यों करता है और SN1 से अभिक्रिया क्यों नहीं करता।
\end{exercise}

\begin{exercise}[$\star\star$]\label{exo:b1:nucleophilic-substitution:7}
\cip{S}-2-आयोडोऑक्टेन को सोडियम आयोडाइड के साथ प्रोपेनोन में रखने पर वह अपनी \omterm{def:b1:stereochemistry-in-depth:optical-activity}{प्रकाशिक सक्रियता} खो देता है,
यद्यपि कोई नया यौगिक नहीं बनता। व्याख्या कीजिए, और बताइए कि
सक्रियता की हानि की दर प्रतिस्थापन की दर की दुगुनी क्यों है।
\end{exercise}

\begin{exercise}[$\star\star$]\label{exo:b1:nucleophilic-substitution:8}
आयोडाइड के साथ SN2 की दर के क्रम में रखिए: 2-ब्रोमो-2-मेथिलप्रोपेन, ब्रोमोमेथेन,
2-ब्रोमोप्रोपेन, 1-ब्रोमोप्रोपेन। इन्हीं यौगिकों को SN1
\omterm{def:b1:nucleophilic-substitution:sn1}{विलायक-अपघटन} की दर के क्रम में रखिए।
\end{exercise}

\begin{exercise}[$\star\star$]\label{exo:b1:nucleophilic-substitution:9}
प्रकाशतः शुद्ध एक द्वितीयक ब्रोमाइड का \omterm{def:b1:nucleophilic-substitution:sn1}{विलायक-अपघटन} होता है; उत्पाद
प्रकाश को शुद्ध प्रतीपन के लिए अपेक्षित मान के \qty{12}{\percent} तक घुमाता है।
उत्पाद के दोनों \omterm{def:b1:stereochemistry-in-depth:enantiomers}{प्रतिबिंबरूपियों} के अनुपात बताइए और
व्याख्या कीजिए।
\end{exercise}

\begin{exercise}[$\star\star\star$]\label{exo:b1:nucleophilic-substitution:10}
\omterm{def:b1:elementary-steps:qssa}{स्थायी-अवस्था सन्निकटन} से सामान्य SN1 \omterm{def:b1:rate-laws:order}{दर नियम} व्युत्पन्न कीजिए, और
दिखाइए कि आरंभ में \omterm{def:b1:electronic-effects:nucleophile}{निर्गामी समूह} आयन \ce{Y-} मिलाने से अभिक्रिया
धीमी हो जाती है (गतिकी का \omterm{def:b1:precipitation:common-ion}{सम-आयन प्रभाव})। किस शर्त पर
यह प्रभाव नगण्य है?
\end{exercise}

\begin{exercise}[$\star\star\star$]\label{exo:b1:nucleophilic-substitution:11}
हाइड्रॉक्साइड के साथ ब्रोमोमेथेन की अभिक्रिया की दो \omterm{def:b1:rate-laws:half-life}{अर्ध-आयु}:
$[\ce{OH-}]_0 = [\ce{CH3Br}]_0 = C_0$ के साथ, दिखाइए कि $1/[\ce{CH3Br}] = 1/C_0 +
kt$ और $t_{1/2}$ बताइए। $[\ce{OH-}]_0 = 50\,C_0$ के साथ, दिखाइए कि
क्षय चरघातांकी है और $t_{1/2}$ बताइए। अभ्यास के आँकड़े: $k =
\qty{2.0e-4}{L/(mol.s)}$, $C_0 = \qty{0.010}{mol/L}$।
\end{exercise}

\begin{exercise}[$\star\star\star$]\label{exo:b1:nucleophilic-substitution:12}
समझाइए कि कोई ऐल्कोहॉल, \ce{ROH}, सोडियम ब्रोमाइड से अभिक्रिया क्यों नहीं करता, पर
सांद्र हाइड्रोब्रोमिक अम्ल से अभिक्रिया करके \ce{RBr} क्यों देता है। एक तृतीयक और एक
प्राथमिक ऐल्कोहॉल के लिए क्रियाविधि लिखिए।
\end{exercise}

\section{समस्या: टर्ट-ब्यूटिल क्लोराइड का विलायक-अपघटन}

\begin{problem}[{सप्ताहांत समस्या --- चालकता आँकड़ों से कोटि और दर स्थिरांक,
अर्ध-आयु, क्रियाविधि और उसका त्रिविम रसायन,
और विलायक-अपघटन की सक्रियण ऊर्जा}]
\label{pb:b1:nucleophilic-substitution:1}
2-क्लोरो-2-मेथिलप्रोपेन (टर्ट-ब्यूटिल क्लोराइड) को कम
सांद्रता पर जल--प्रोपेनोन मिश्रण में घोला जाता है। उसका \omterm{def:b1:nucleophilic-substitution:sn1}{विलायक-अपघटन}
\ce{(CH3)3CCl + H2O -> (CH3)3COH + H+ + Cl-} आयन उत्पन्न करता है, और
विलयन की चालकता $\kappa$, जो बने आयनों की मात्रा के
समानुपाती है, अभिलिखित की जाती है। बहुत लंबे समय के बाद वह दोनों तापों पर
$\kappa_\infty = \qty{2.000}{mS/cm}$ तक पहुँचती है। मापन
(\unit{mS/cm}):

\begin{center}
\begin{tabular}{lcccccccc}
\hline
$t$ (मिनट) & 0 & 5 & 10 & 15 & 20 & 30 & 45 & 60 \\
\hline
$\kappa$, \qty{25.0}{\celsius} & 0 & 0.172 & 0.329 & 0.473 & 0.605 & 0.835 & 1.110 & 1.321 \\
$\kappa$, \qty{35.0}{\celsius} & 0 & 0.507 & 0.885 & 1.168 & 1.379 & 1.654 & 1.856 & 1.940 \\
\hline
\end{tabular}
\end{center}
% data generated in figdata/bachelor-1/solvolysis.py (story data)

\medskip
\noindent\textbf{भाग I --- कोटि।}
\begin{enumerate}
  \item चालकता अभिक्रिया कर चुके \omterm{def:b1:nucleophilic-substitution:substitution}{क्रियाधार} की मात्रा के समानुपाती
        क्यों है?
  \item $[\mathrm{RCl}]/[\mathrm{RCl}]_0$ को $\kappa$ और
        $\kappa_\infty$ के पदों में व्यक्त कीजिए।
  \item प्रथम कोटि की अभिक्रिया का समाकलित नियम लिखिए और वह
        राशि बताइए जिसे इसकी जाँच के लिए $t$ के विरुद्ध आलेखित करना है।
  \item हर समय के लिए \qty{25.0}{\celsius} पर $\ln(\kappa_\infty - \kappa)$
        परिकलित कीजिए।
  \item क्या आलेख सरल रेखा है? कोटि के बारे में निष्कर्ष निकालिए।
  \item जल बड़े आधिक्य में है। यह कोटि में क्या छिपा देता है?
\end{enumerate}

\medskip
\noindent\textbf{भाग II --- \omterm{def:b1:rate-laws:order}{दर स्थिरांक} और \omterm{def:b1:rate-laws:half-life}{अर्ध-आयु}।}
\begin{enumerate}[resume]
  \item \qty{25.0}{\celsius} पर \omterm{def:b1:rate-laws:order}{दर स्थिरांक} \unit{s^{-1}} में ज्ञात कीजिए।
  \item \omterm{def:b1:rate-laws:half-life}{अर्ध-आयु} परिकलित कीजिए।
  \item किस समय पर \qty{90}{\percent} \omterm{def:b1:nucleophilic-substitution:substitution}{क्रियाधार} खप जाता है?
  \item \qty{35.0}{\celsius} पर \omterm{def:b1:rate-laws:order}{दर स्थिरांक} ज्ञात कीजिए।
  \item दस डिग्री के लिए दर कितने गुना बढ़ी है?
  \item \qty{35}{\celsius} शृंखला के एक मान की नियम से जाँच कीजिए।
\end{enumerate}

\medskip
\noindent\textbf{भाग III --- क्रियाविधि।}
\begin{enumerate}[resume]
  \item \omterm{def:b1:nucleophilic-substitution:substitution}{क्रियाधार} और विलायक कौन-सी क्रियाविधि सुझाते हैं?
  \item उसे \omterm{def:b1:electronic-effects:curly-arrow}{वक्र तीरों} के साथ लिखिए।
  \item दिखाइए कि यह प्राप्त कोटि से सहमत है।
  \item \qty{0.01}{mol/L} पर सोडियम हाइड्रॉक्साइड मिलाने से क्लोराइड के \omterm{def:b1:rate-laws:rate}{लोप की
        दर} में शायद ही कोई परिवर्तन होता है। क्यों?
  \item \omterm{def:b1:stereochemistry-in-depth:stereocentre}{किरैल} तृतीयक क्लोराइड
        3-क्लोरो-3-मेथिलहेक्सेन के साथ, प्रकाशतः शुद्ध, वही प्रयोग लगभग
        शून्य \omterm{def:b1:stereochemistry-in-depth:optical-activity}{प्रकाशिक सक्रियता} वाला ऐल्कोहॉल देता है। व्याख्या कीजिए।
  \item अभिक्रिया का \omterm{def:b1:elementary-steps:energy-profile}{ऊर्जा प्रोफ़ाइल} खींचिए और
        \omterm{def:b1:elementary-steps:rds}{दर-निर्धारक चरण} चिह्नित कीजिए।
  \item उन्हीं परिस्थितियों में 1-क्लोरोब्यूटेन शायद ही अभिक्रिया क्यों करेगा?
\end{enumerate}

\medskip
\noindent\textbf{भाग IV --- \omterm{def:b1:rate-laws:activation-energy}{सक्रियण ऊर्जा}।}
\begin{enumerate}[resume]
  \item \omterm{thm:b1:rate-laws:arrhenius}{आरेनियस नियम} लिखिए।
  \item दोनों \omterm{def:b1:rate-laws:order}{दर स्थिरांकों} से $E_a$ व्यक्त कीजिए।
  \item $E_a$ परिकलित कीजिए ($R = \qty{8.314}{J/(K.mol)}$)।
  \item क्रियाविधि में यह ऊर्जा क्या मापती है?
  \item \qty{45.0}{\celsius} पर \omterm{def:b1:rate-laws:order}{दर स्थिरांक} का पूर्वानुमान कीजिए।
  \item \omterm{def:b1:nucleophilic-substitution:sn1}{विलायक-अपघटन} की \omterm{def:b1:rate-laws:activation-energy}{सक्रियण ऊर्जा} दो सार्थक
        अंकों तक बताइए।
\end{enumerate}
\end{problem}
